\documentclass[11pt]{amsart}

\usepackage{amsmath,amssymb,amsthm,mathtools}
\usepackage{enumitem}
\usepackage[colorlinks=true,citecolor=blue,linkcolor=blue,urlcolor=blue]{hyperref}

\newtheorem{theorem}{Theorem}[section]
\newtheorem{proposition}[theorem]{Proposition}
\newtheorem{lemma}[theorem]{Lemma}
\newtheorem{corollary}[theorem]{Corollary}

\theoremstyle{definition}
\newtheorem{definition}[theorem]{Definition}
\newtheorem{remark}[theorem]{Remark}

\newcommand{\V}{\mathcal V}
\newcommand{\E}{\mathcal E}
\newcommand{\Lk}{\operatorname{Lk}}
\newcommand{\St}{\operatorname{St}}
\newcommand{\Tr}{\operatorname{Tr}}
\newcommand{\B}{\mathbb B}
\newcommand{\K}{\mathbb K}
\newcommand{\C}{\mathbb C}
\newcommand{\bbarotimes}{\mathbin{\overline\otimes}}
\newcommand{\freeprod}{\mathop{*}}
\newcommand{\Ecal}{\mathcal E}

\begin{document}
\title{Submission C solution to Problem 10}
\date{}
\maketitle



\begin{abstract}
We prove that if $\Gamma$ is a finite join-irreducible simplicial graph with at least three vertices and if each tracial vertex algebra $(M_v,\tau_v)$ contains a trace-zero unitary, then the tracial graph product von Neumann algebra is properly proximal.
\end{abstract}

\section{Preliminaries}

We use the following standard convention.  A finite graph $\Gamma=(\V,\E)$ is called \emph{irreducible} if it is not a non-trivial join: equivalently, there is no non-empty proper subset $S\subset \V$ such that every vertex of $S$ is adjacent to every vertex of $\V\setminus S$.  Equivalently, the complement graph $\Gamma^c$ is connected.

For $U\subset \V$, let $\Gamma_U$ be the induced subgraph and let
\[
(M_U,\tau_U)=\mathop{\star}_{v\in \Gamma_U}(M_v,\tau_v)
\]
denote the corresponding graph product von Neumann algebra.  We put $M_\varnothing=\C$.  The canonical trace-preserving conditional expectation is denoted
\[
E_U:M\to M_U.
\]
We write
\[
\Lk(v)=\{w\in \V: (v,w)\in \E\},
\qquad
\St(v)=\Lk(v)\cup\{v\}.
\]

We use the definition of proper proximality for finite von Neumann algebras introduced by Ding--Kunnawalkam Elayavalli--Peterson \cite{DKEP}.  If $N\subset M$ is a von Neumann subalgebra with expectation, let $X_N$ denote the $M$-boundary piece associated with $N$, i.e. the boundary piece generated by the Jones projection $e_N:L^2(M)\to L^2(N)$.  The small-at-infinity operator system is denoted $S_{X_N}(M)$.  By \cite[Theorems 6.2 and 8.5]{DKEP}, $M$ is properly proximal relative to $X_N$ exactly when there is no $M$-central state on $S_{X_N}(M)$ whose restriction to $M$ is normal.

We shall use two standard consequences of the theory in \cite{DKEP}.

\begin{lemma}[Upgrading from relative boundary pieces]\label{lem:upgrade}
Let $(M,\tau)$ be separable and tracial.  Let $X_1,\dots,X_n$ be $M$-boundary pieces and let $q_i$ be the support projection of $K_{X_i}(M)$ in the bidual model of \cite[Section 8]{DKEP}.  Let $q_{\K}$ be the corresponding support projection for the compact boundary piece $\K(L^2M)$.  Suppose:
\begin{enumerate}[label=(\roman*)]
\item $M$ is properly proximal relative to every $X_i$;
\item $\bigwedge_{i=1}^n q_i=q_{\K}$.
\end{enumerate}
Then $M$ is properly proximal.
\end{lemma}

\begin{proof}
If $M$ were not properly proximal, then \cite[Theorem 8.5]{DKEP} would give an $M$-central state $\varphi$ on the bidual small-at-infinity system for the compact boundary whose restriction to $M$ is normal.  Since proper proximality relative to any $X_i$ rules out amenable direct summands, $\varphi(q_{\K})=0$; otherwise Connes' amenability criterion would produce an amenable direct summand.

For each $i$, if $\varphi(q_i^\perp)>0$, then the normalized compression of $\varphi$ by $q_i^\perp$ gives an $M$-central state on the small-at-infinity system relative to $X_i$, still normal on $M$, contradicting relative proper proximality.  Hence $\varphi(q_i)=1$ for all $i$.  Since the $q_i$ are central projections and the family is finite,
\[
\varphi\Big(\bigwedge_i q_i\Big)=1.
\]
By assumption, $\bigwedge_i q_i=q_{\K}$, contradicting $\varphi(q_{\K})=0$.
\end{proof}

\begin{lemma}[Boundary cover for graph products]\label{lem:boundary-cover}
For a graph product $M=\star_{\Gamma}(M_v,\tau_v)$ and subsets $U_1,\dots,U_n\subset \V$, let $q_{U_i}$ be the support projection associated with the boundary piece $X_{M_{U_i}}$.  Then
\[
q_{U_1}\wedge\cdots\wedge q_{U_n}=q_{U_1\cap\cdots\cap U_n}.
\]
In particular,
\[
\bigwedge_{v\in \V} q_{\Lk(v)}=q_\varnothing=q_{\K}.
\]
\end{lemma}

\begin{proof}
The first assertion is the support-projection form of the graph-product bimodule calculation for subgraph algebras; see Charlesworth--de Santiago--Hayes--Jekel--Kunnawalkam Elayavalli--Nelson \cite[Theorem 0.4, Theorem 5.4, and Proposition 5.5]{CDHJKEN}.  Equivalently, the $M$-$M$ boundary piece generated by $e_{M_U}$ intersects/fuses with the one generated by $e_{M_W}$ exactly as the subgraph algebras intersect:
\[
M_U\cap M_W=M_{U\cap W},
\qquad
E_U E_W=E_{U\cap W}.
\]
Iterating gives the formula for finitely many $U_i$.

For $U_i=\Lk(v_i)$ over all vertices, the intersection is empty because $v\notin \Lk(v)$ for every $v$.  Thus the meet is $q_\varnothing$.  Since $M_\varnothing=\C$ and the boundary piece generated by $e_{\C}$ is exactly $\K(L^2M)$, we have $q_\varnothing=q_{\K}$.
\end{proof}

\section{A relative proper proximality criterion for amalgamated free products}

We need the following elementary relative version of the usual paradoxical-decomposition proof for free products.

\begin{proposition}\label{prop:AFP}
Let
\[
(M,\tau)=(A,E_B^A)\freeprod_B(C,E_B^C)
\]
be a tracial amalgamated free product over a common von Neumann subalgebra $B$.  Suppose there are unitaries
\[
u_1,u_2\in C,\qquad w\in A
\]
such that
\[
E_B^C(u_1)=E_B^C(u_2)=E_B^C(u_1^*u_2)=0,
\qquad
E_B^A(w)=0.
\]
Then $M$ is properly proximal relative to the boundary piece $X_B$ associated with $B$.
\end{proposition}

\begin{proof}
Let $A^\circ=\ker E_B^A$ and $C^\circ=\ker E_B^C$.  In the standard $L^2$-model of the amalgamated free product,
\[
L^2(M)=L^2(B)\oplus H_A\oplus H_C,
\]
where $H_A$ is the closed span of reduced words beginning in $A^\circ$, and $H_C$ is the closed span of reduced words beginning in $C^\circ$.  Let $P_A$ and $P_C$ be the corresponding orthogonal projections.

A direct Fock-space computation shows that
\[
P_A,P_C\in S_{X_B}(M).
\]
Indeed, right multiplication by an algebraic reduced word changes the first letter only through cancellations that factor through $L^2(B)$; hence the commutators with the right action of $M$ lie in the ideal generated by $e_B$, namely $K_{X_B}(M)$.  This is the standard computation used in \cite[Section 6]{DKEP}.

Assume, toward contradiction, that there exists an $M$-central state $\varphi$ on $S_{X_B}(M)$ whose restriction to $M$ is normal.  Since $u_1,u_2\in C^\circ$ and $E_B(u_1^*u_2)=0$, the subspaces $u_1^*H_A$ and $u_2^*H_A$ are orthogonal and both are contained in $H_C$.  Hence
\[
u_1^*P_Au_1+u_2^*P_Au_2\leq P_C.
\]
By $M$-centrality and unitarity,
\[
2\varphi(P_A)
=
\varphi(u_1^*P_Au_1)+\varphi(u_2^*P_Au_2)
\leq \varphi(P_C).
\]
Similarly, since $w\in A^\circ$,
\[
w^*P_Cw\leq P_A,
\]
and therefore
\[
\varphi(P_C)=\varphi(w^*P_Cw)\leq \varphi(P_A).
\]
Thus $\varphi(P_A)=\varphi(P_C)=0$.

Finally, since $u_1\in C^\circ$, the range of $u_1^*e_Bu_1$ is contained in $H_C$, so
\[
u_1^*e_Bu_1\leq P_C.
\]
Again by centrality,
\[
\varphi(e_B)=\varphi(u_1^*e_Bu_1)\leq \varphi(P_C)=0.
\]
But $1=P_A+P_C+e_B$, so $\varphi(1)=0$, a contradiction.
\end{proof}

\section{Graph-product input}

For every vertex $v$, the graph product admits the canonical amalgamated free product decomposition
\[
M
=
M_{\St(v)}
\freeprod_{M_{\Lk(v)}}
M_{\V\setminus\{v\}}.
\]
This follows directly from the graph-product normal form and is standard; see \cite{CF} and \cite[Section 1.1]{CDHJKEN}.

We now verify that the hypotheses of Proposition \ref{prop:AFP} hold for each vertex.

\begin{lemma}\label{lem:unitaries}
Let $\Gamma$ be finite, irreducible, and $|\V|\geq 3$.  Fix $v\in\V$ and put $B=M_{\Lk(v)}$.  In the decomposition
\[
M
=
M_{\St(v)}
\freeprod_B
M_{\V\setminus\{v\}},
\]
there exist a unitary $w\in M_{\St(v)}$ and unitaries $u_1,u_2\in M_{\V\setminus\{v\}}$ such that
\[
E_B(w)=0,\qquad
E_B(u_1)=E_B(u_2)=E_B(u_1^*u_2)=0.
\]
\end{lemma}

\begin{proof}
Choose, for each vertex $s$, a trace-zero unitary $a_s\in M_s$.

First take
\[
w=a_v\in M_v\subset M_{\St(v)}.
\]
Since $M_{\St(v)}=M_v\bbarotimes M_{\Lk(v)}$ over the commuting graph product decomposition, we have
\[
E_B(w)=\tau_v(a_v)1=0.
\]

Let
\[
L=\Lk(v),\qquad W=\V\setminus(\{v\}\cup L).
\]
The set $W$ is non-empty.  Indeed, if $W=\varnothing$, then $v$ is adjacent to every vertex in $\V\setminus\{v\}$, giving the non-trivial join decomposition
\[
\Gamma=\{v\}+\Gamma_{\V\setminus\{v\}},
\]
contrary to irreducibility.

We distinguish two cases.

\smallskip

\noindent\textbf{Case 1: $L=\varnothing$.}
Then $B=\C$.  Since $|\V|\geq 3$, choose distinct vertices $r,s\in \V\setminus\{v\}$.  Put
\[
u_1=a_r,\qquad u_2=a_s a_r.
\]
Both are unitaries in $M_{\V\setminus\{v\}}$.  Their scalar expectations are zero:
\[
\tau(u_1)=0,\qquad \tau(u_2)=0.
\]
Moreover,
\[
u_1^*u_2=a_r^*a_s a_r.
\]
If $r$ and $s$ are adjacent, this equals $a_s$, whose trace is zero.  If they are not adjacent, it is a non-empty reduced word with trace-zero letters, hence has trace zero.  Thus $E_B(u_1^*u_2)=0$.

\smallskip

\noindent\textbf{Case 2: $L\neq \varnothing$.}
We claim that there exist $\ell\in L$ and $r\in W$ such that $\ell$ is not adjacent to $r$.  If not, then every vertex of $L$ is adjacent to every vertex of $\V\setminus L$ because vertices of $L$ are adjacent to $v$ by definition and, by assumption, to every vertex of $W$.  Hence
\[
\Gamma=\Gamma_L+\Gamma_{\V\setminus L},
\]
again contradicting irreducibility.

Choose such $\ell\in L$ and $r\in W$.  Put
\[
u_1=a_r,\qquad u_2=a_\ell a_r.
\]
Since $a_\ell\in M_\ell\subset B$ and $r\notin L$, both $u_1$ and $u_2$ lie in $M_{\V\setminus\{v\}}$ and satisfy
\[
E_B(u_1)=E_B(u_2)=0.
\]
Furthermore,
\[
u_1^*u_2=a_r^*a_\ell a_r.
\]
Because $r$ and $\ell$ are not adjacent, this is a reduced word involving the vertex $r\notin L$, so its expectation onto $B=M_L$ is zero.  Hence $E_B(u_1^*u_2)=0$.
\end{proof}

\begin{corollary}\label{cor:relative}
For every $v\in\V$, the graph product $M$ is properly proximal relative to the boundary piece $X_{M_{\Lk(v)}}$.
\end{corollary}

\begin{proof}
Apply Proposition \ref{prop:AFP} to
\[
M
=
M_{\St(v)}
\freeprod_{M_{\Lk(v)}}
M_{\V\setminus\{v\}},
\]
using the unitaries supplied by Lemma \ref{lem:unitaries}.
\end{proof}

\section{Proof of the theorem}

\begin{theorem}
Let $\Gamma$ be a finite simplicial graph with $|\V|\geq 3$.  Let $(M_v,\tau_v)$ be a separable tracial von Neumann algebra for every $v\in \V$.  Assume:
\begin{enumerate}[label=(\alph*)]
\item $\Gamma$ is irreducible, i.e. not a non-trivial join;
\item every $M_v$ contains a trace-zero unitary.
\end{enumerate}
Let
\[
(M,\tau)=\mathop{\star}_{v\in\Gamma}(M_v,\tau_v)
\]
be the tracial graph product von Neumann algebra.  Then $M$ is properly proximal.
\end{theorem}

\begin{proof}
For every vertex $v$, Corollary \ref{cor:relative} shows that $M$ is properly proximal relative to the boundary piece $X_{M_{\Lk(v)}}$.

Therefore $M$ has no amenable direct summand: if $pM$ were an amenable direct summand, Connes' amenability criterion would give an $M$-central state on $\B(L^2(pM))$, hence by extension an $M$-central state on every $S_{X_{M_{\Lk(v)}}}(M)$ normal on $M$, contradicting relative proper proximality.

By Lemma \ref{lem:boundary-cover},
\[
\bigwedge_{v\in\V} q_{\Lk(v)}=q_{\K}.
\]
Thus the finite family of boundary pieces $\{X_{M_{\Lk(v)}}:v\in\V\}$ satisfies the hypotheses of the upgrading Lemma \ref{lem:upgrade}.  Consequently $M$ is properly proximal.
\end{proof}

\begin{thebibliography}{99}

\bibitem{CF}
M.~Caspers and P.~Fima,
\emph{Graph products of operator algebras},
J. Noncommut. Geom. \textbf{11} (2017), no.~1, 367--411.

\bibitem{CDHJKEN}
I.~Charlesworth, R.~de Santiago, B.~Hayes, D.~Jekel,
S.~Kunnawalkam Elayavalli and B.~Nelson,
\emph{On the structure of graph product von Neumann algebras},
Publ. Res. Inst. Math. Sci. \textbf{61} (2025), no.~4, 713--762.

\bibitem{DKEP}
C.~Ding, S.~Kunnawalkam Elayavalli and J.~Peterson,
\emph{Properly proximal von Neumann algebras},
Duke Math. J. \textbf{172} (2023), no.~15, 2821--2894.

\end{thebibliography}

\end{document}